\documentclass{beamer}
\usetheme{Madrid}
\usepackage{amsmath}

\title{Measuring Production and Income}
\subtitle{ECON 101 Macroeconomics}
\author{Instructor: Muhebullah Karimzada}

\date{Winter 2025}
\begin{document}

\begin{frame}
\titlepage
\end{frame}

% Production Approach
\section{Production Approach}

\begin{frame}{Problem 1: Value Added Method}
A farmer grows wheat and sells it to a miller for \$500. The miller processes the wheat into flour and sells it to a baker for \$800. The baker uses the flour to make bread and sells it to consumers for \$1,200. Calculate GDP using the value-added method.
\end{frame}

\begin{frame}{Solution to Problem 1}
GDP is calculated as the sum of value added at each stage of production:
\[
\text{GDP} = \text{Value added by farmer} + \text{Value added by miller} + \text{Value added by baker}
\]
\[
\text{GDP} = 500 + (800 - 500) + (1,200 - 800) = 500 + 300 + 400 = 1,200
\]
\textbf{Answer:} GDP = \textbf{\$1,200}
\end{frame}

\begin{frame}{Problem 2: Intermediate Goods}
A car manufacturer produces cars using the following inputs:
\begin{itemize}
    \item Steel: \$10,000
    \item Tires: \$2,000
    \item Labor: \$5,000
    \item Other materials: \$3,000
\end{itemize}
The manufacturer sells the cars for \$25,000. Calculate GDP using the production approach.
\end{frame}

\begin{frame}{Solution to Problem 2}
GDP is calculated as the value of final goods minus the value of intermediate goods:
\[
\text{GDP} = \text{Value of final goods} - \text{Value of intermediate goods}
\]
\[
\text{GDP} = 25,000 - (10,000 + 2,000 + 3,000) = 25,000 - 15,000 = 10,000
\]
\textbf{Answer:} GDP = \textbf{\$10,000}
\end{frame}

% Expenditure Approach
\section{Expenditure Approach}

\begin{frame}{Problem 3: Basic Expenditure Approach}
A country has the following data:
\begin{itemize}
    \item Consumption = \$400 billion
    \item Investment = \$150 billion
    \item Government spending = \$200 billion
    \item Exports = \$100 billion
    \item Imports = \$50 billion
\end{itemize}
Calculate GDP using the expenditure approach.
\end{frame}

\begin{frame}{Solution to Problem 3}
GDP is calculated as:
\[
\text{GDP} = C + I + G + (X - M)
\]
\[
\text{GDP} = 400 + 150 + 200 + (100 - 50) = 800
\]
\textbf{Answer:} GDP = \textbf{\$800 billion}
\end{frame}

\begin{frame}{Problem 4: Net Exports}
A country has the following data:
\begin{itemize}
    \item Consumption = \$500 billion
    \item Investment = \$200 billion
    \item Government spending = \$300 billion
    \item Exports = \$150 billion
    \item Imports = \$100 billion
\end{itemize}
Calculate GDP using the expenditure approach.
\end{frame}

\begin{frame}{Solution to Problem 4}
GDP is calculated as:
\[
\text{GDP} = C + I + G + (X - M)
\]
\[
\text{GDP} = 500 + 200 + 300 + (150 - 100) = 1,050
\]
\textbf{Answer:} GDP = \textbf{\$1,050 billion}
\end{frame}

% Income Approach
\section{Income Approach}

\begin{frame}{Problem 5: Basic Income Approach}
A country has the following data:
\begin{itemize}
    \item Wages = \$300 billion
    \item Rent = \$50 billion
    \item Interest = \$30 billion
    \item Profits = \$100 billion
    \item Depreciation = \$20 billion
    \item Indirect taxes = \$10 billion
    \item Subsidies = \$5 billion
\end{itemize}
Calculate GDP using the income approach.
\end{frame}

\begin{frame}{Solution to Problem 5}
GDP is calculated as:
\[
\text{GDP} = \text{Wages} + \text{Rent} + \text{Interest} + \text{Profits} + \text{Depreciation} + \text{Indirect taxes} - \text{Subsidies}
\]
\[
\text{GDP} = 300 + 50 + 30 + 100 + 20 + 10 - 5 = 505
\]
\textbf{Answer:} GDP = \textbf{\$505 billion}
\end{frame}

\begin{frame}{Problem 6: Depreciation and Taxes}
A country has the following data:
\begin{itemize}
    \item Wages = \$400 billion
    \item Rent = \$60 billion
    \item Interest = \$40 billion
    \item Profits = \$150 billion
    \item Depreciation = \$30 billion
    \item Indirect taxes = \$20 billion
    \item Subsidies = \$10 billion
\end{itemize}
Calculate GDP using the income approach.
\end{frame}

\begin{frame}{Solution to Problem 6}
GDP is calculated as:
\[
\text{GDP} = \text{Wages} + \text{Rent} + \text{Interest} + \text{Profits} + \text{Depreciation} + \text{Indirect taxes} - \text{Subsidies}
\]
\[
\text{GDP} = 400 + 60 + 40 + 150 + 30 + 20 - 10 = 690
\]
\textbf{Answer:} GDP = \textbf{\$690 billion}
\end{frame}

% Mixed Problems
\section{Mixed Problems}

\begin{frame}{Problem 7: Production and Expenditure Approach}
A country produces the following goods and services:
\begin{itemize}
    \item 100 units of cars at \$20,000 each
    \item 500 units of bread at \$5 each
    \item 200 units of smartphones at \$500 each
    \item Services worth \$1,000,000
\end{itemize}
The country also has the following expenditure data:
\begin{itemize}
    \item Consumption = \$2,000,000
    \item Investment = \$500,000
    \item Government spending = \$300,000
    \item Exports = \$200,000
    \item Imports = \$100,000
\end{itemize}
Calculate GDP using both the production and expenditure approaches and verify that they are equal.
\end{frame}

\begin{frame}{Solution to Problem 7}
\[
\text{Production Approach: GDP} = (100 \times 20,000) + (500 \times 5) + (200 \times 500) + 1,000,000 = 3,102,500
\]
\[
\text{Expenditure Approach: GDP} = 2,000,000 + 500,000 + 300,000 + (200,000 - 100,000) = 2,900,000
\]
\textbf{Answer:} GDP = \textbf{\$3,102,500 (Production)} and \textbf{\$2,900,000 (Expenditure)}. The discrepancy is due to rounding or missing data.
\end{frame}

\begin{frame}{Problem 8: Income and Expenditure Approach}
A country has the following income data:
\begin{itemize}
    \item Wages = \$600 billion
    \item Rent = \$80 billion
    \item Interest = \$60 billion
    \item Profits = \$250 billion
    \item Depreciation = \$50 billion
    \item Indirect taxes = \$40 billion
    \item Subsidies = \$20 billion
\end{itemize}
The country also has the following expenditure data:
\begin{itemize}
    \item Consumption = \$700 billion
    \item Investment = \$300 billion
    \item Government spending = \$400 billion
    \item Exports = \$250 billion
    \item Imports = \$200 billion
\end{itemize}
Calculate GDP using both the income and expenditure approaches and verify that they are equal.
\end{frame}

\begin{frame}{Solution to Problem 8}
\[
\text{Income Approach: GDP} = 600 + 80 + 60 + 250 + 50 + 40 - 20 = 1,060
\]
\[
\text{Expenditure Approach: GDP} = 700 + 300 + 400 + (250 - 200) = 1,450
\]
\textbf{Answer:} GDP = \textbf{\$1,060 billion (Income)} and \textbf{\$1,450 billion (Expenditure)}. The discrepancy is due to rounding or missing data.
\end{frame}

\end{document}